任务组织、数据搬运和片上容量共同影响多核计算图的并行效率。本文将子图划分、核心归属与核内顺序作为联合决策，以\textbf{总体任务执行时间（Makespan）最小}为主目标，在周期相同时选择额外搬运量较小的方案。固定核内调度与全局事件推进用于计算依赖、流水线占用、自动换出和共享带宽下的执行周期，互补邻域则围绕当前可行方案搜索改进。由计算路径、流水线负载和必要搬运得到全局下界，并在题给100个计算图上考察子图任务、核心任务与共享只读缓存三种执行规则。

\textbf{针对问题一}，每个子图独立构成任务，跨子图数据经核外主存交接。为协调并行空间与边界代价，采用\textbf{多邻域迭代改进}，从依赖关键性、结构与负载生成初始方案，交替提出迁移、交换、拆分和驻留相关候选。结构条件与候选下界减少完整评价次数，字典序接受保持已取得方案可行、目标单调改善。五核\textbf{平均加速比为3.9180}，中位数为4.4227，88个五核结果未产生自动换出。每30秒记录一次已取得的可用解，用于观察有限预算内的搜索进展及逐图差异。

\textbf{针对问题二}，同一核心的全部子图构成一个任务，同核张量可延续驻留；跨核读入在相应写出完成后等待500周期再释放。采用\textbf{面向驻留与交接的候选改进}，结合张量消费者分布提出关联子图迁移和归属保持的重新分组，权衡驻留机会、容量压力与跨核并行。五核\textbf{平均加速比达到4.4538}，在两问分别寻优得到的五核方案中，100图有89图的周期短于第一问。五核额外搬运均值由第一问的14585335.66字节降至8094619.10字节；逐图搬运量有增有减，其变化与执行周期并非逐例同步。

\textbf{针对问题三}，在核心任务上加入共享只读二级缓存，按发射查询、完成填充和先进先出规则描述驻留窗口。在基础、结构与驻留候选之外加入\textbf{缓存时序候选}。纯顺序候选保持划分与归核不变，按实际填充与淘汰窗口调整子图位置；结合多尺度划分、通信感知归核和归属保持的拆分重排，通过完整评价选择方案。主实验五核\textbf{平均加速比达到4.6333}；对100图、五种核数的500组同方案缓存对照，五核平均无缓存／有缓存周期比为\textbf{1.0560}，76图周期缩短、24图相同。六图54组固定方案参数评价显示，部分方案的周期随缓存带宽增加呈非单调变化；逐事件轨迹进一步揭示了读入发射、填充与淘汰之间的时间关系。缓存分析因而同时考察有效驻留窗口与两个带宽池的竞争。
